\documentclass[english,onecolumn]{IEEEtran}
\usepackage[T1]{fontenc}
\usepackage[colorlinks]{hyperref}
\usepackage{color,xcolor}
\usepackage{amsthm,amssymb,amsfonts,amsmath}
\usepackage{mathtools}
\usepackage{algorithm}
\usepackage{algpseudocode}
% for matlab code highlight
% load package with ``framed'' and ``numbered'' option.
\usepackage[framed,numbered,autolinebreaks,useliterate]{mcode}

\providecommand{\U}[1]{\protect\rule{.1in}{.1in}}
\topmargin            -18.0mm
\textheight           226.0mm
\oddsidemargin      -4.0mm
\textwidth            166.0mm
\def\baselinestretch{1.5}




% math operator macros
\newcommand{\madj}{\mathop{\mathrm{adj}}}
\newcommand{\trace}{\mathop{\mathrm{tr}}}


\DeclarePairedDelimiter{\abs}{\lvert}{\rvert}
\DeclarePairedDelimiter{\norm}{\lVert}{\rVert} 
\DeclarePairedDelimiter{\pnorm}{\lVert}{\rVert_p}



% special letters
\newcommand{\x}{{\mathbf{x}}} % vector x
\newcommand{\y}{{\mathbf{y}}} % vector y
\newcommand{\z}{{\mathbf{z}}} % vector z
\newcommand{\0}{{\mathbf{0}}} % zero vector

\newcommand{\A}{{\mathbf{A}}} % matrix A
\newcommand{\B}{{\mathbf{B}}} % matrix B
\newcommand{\matC}{{\mathbf{C}}} % matrix C
\newcommand{\matD}{{\mathbf{D}}} % matrix D
\newcommand{\matE}{{\mathbf{E}}} % matrix E
\newcommand{\matH}{{\mathbf{H}}} % matrix H
\newcommand{\matK}{{\mathbf{K}}} % matrix K
\newcommand{\X}{{\mathbf{X}}} % matrix X
\newcommand{\Y}{{\mathbf{Y}}} % matrix Y
\newcommand{\Z}{{\mathbf{Z}}} % matrix Z
\newcommand{\M}{{\mathbf{M}}} % matrix M
\newcommand{\I}{{\mathbf{I}}} % identity matrix

\newcommand{\calV}{{\mathcal{V}}} % subspace V
\newcommand{\calU}{{\mathcal{U}}} % subspace U
\newcommand{\calS}{{\mathcal{S}}} % subspace S
\newcommand{\calM}{{\mathcal{M}}} % subspace M
\newcommand{\calN}{{\mathcal{N}}} % subspace N
\newcommand{\calR}{\mathop{\mathcal{R}}} % range
\renewcommand{\calN}{\mathop{\mathcal{N}}} % nullspace

% number fields
\newcommand{\R}{\mathbb{R}} % real numbers
\newcommand{\C}{\mathbb{C}} % complex numbers

\newcommand{\bigO}{{\mathcal{O}}} % big-Oh notation









\begin{document}

\begin{center}
	\textbf{\LARGE{SI231B: Matrix Computations, 2025 Fall}}\\
	{\Large Homework Set \#4}\\
	% \texttt{Prof. Jie Lu}
\par\end{center}

\noindent
\rule{\linewidth}{0.4pt}
% \noindent
% \rule{\linewidth}{0.4pt}
{\bf Acknowledgements:}
\begin{enumerate}
	\item Deadline: {\bf \textcolor{red}{2025-12-13 23:59:59}}
    \item Please submit the PDF file to \hyperlink{https://www.gradescope.com/}{gradescope}. Course entry code: N2382J.
    \item You have 5 ``free days'' in total for all late homework submissions.
    \item If your homework is handwritten, please make it clear and legible.
    \item All your answers are required to be in English. 
    \item Write down the major steps for deriving the solution; otherwise you may loss points.
\end{enumerate}
\rule{\linewidth}{0.4pt}


% =======================================================================================
\newpage
\noindent \textbf{Problem 1. (Quadratic Form )} (\textcolor{blue}{20 points})

\begin{enumerate}
Consider the following function:
\[
f(\mathbf{x}) = 2x_1^2 + 2x_2^2 + ax_3^2 + 2x_1x_2+2x_1x_3+2x_2x_3
\]


    \item Find a symmetric matrix $\mathbf{A} \in \mathbb{R}^{3 \times 3}$ such that $f(\mathbf{x}) = \mathbf{x}^T \mathbf{A} \mathbf{x}$. (\textcolor{blue}{5 points})
    
    \item Find the condition of $a$ such that $\mathbf{A}$ is PSD. (\textcolor{blue}{5 points})
    
    \item If $a=2$, find the PSD square root of $\mathbf{A}$. (\textcolor{blue}{10 points})\\(Hint: You might need to use Gram-Schmidt to find orthonormal eigenvectors.)

\end{enumerate}

















\bigskip

\noindent{\bf Solution:}






% =======================================================================================
\newpage
\noindent \textbf{Problem 2. (Positive Semidefinite Matrices)} (\textcolor{blue}{20 points}) \\
\noindent Let \( \boldsymbol{A} = \boldsymbol{I} + \boldsymbol{\alpha}\boldsymbol{\alpha}^T \) be an \( n \times n\) real symmetric matrix, where  \( \boldsymbol{\alpha} = [a_1, a_2, \dots, a_n]^T \) is a non-zero real vector, and denote \( \beta = \boldsymbol{\alpha}^T\boldsymbol{\alpha} > 0 \). 

\begin{enumerate}
\item Prove that \( \boldsymbol{A} \) is a symmetric positive definite matrix.
  (\textcolor{blue}{5 points}) 
\item Find all eigenvalues of \( \boldsymbol{A} \) and the corresponding linearly independent eigenvectors (specify the algebraic multiplicity of each eigenvalue).(\textcolor{blue}{5 points}) 
\item Let \( \boldsymbol{C} = (\boldsymbol{A} - \boldsymbol{I})^3\). Find the positive definite square root \( \boldsymbol{C}^{1/2} \) of \( \boldsymbol{C} \).
(\textcolor{blue}{10 points})

\end{enumerate}

\bigskip

\noindent \textbf{Solution:}




% =======================================================================================
\newpage
\noindent \textbf{Problem 3. (Application of Schur Complement)} \hfill (\textcolor{blue}{10 points}) 

\begin{enumerate}
Suppose \(\boldsymbol{X}\) is the \(6 \times 6\) matrix:
\[
\boldsymbol{X} = 
\begin{bmatrix}
2 & 2 & 1 & 0 & 0 & -1  \\
2 & a & 2 & 0 & 0 & 2 \\
1 & 2 & 2 & 0 & 0 & 0  \\
0 & 0 & 0 & 2 & 0 & 0 \\
0 & 0 & 0 & 0 & 2 & 0  \\
-1 & 2 & 0 & 0 & 0 & 2 \\

\end{bmatrix}
\]
Find the condition of \(a\) such that \(\boldsymbol{X}\) is PSD.
\textcolor{blue}{(10 points)}
\end{enumerate}

\bigskip

\noindent \textbf{Solution:}





% =======================================================================================
\newpage
\noindent \textbf{Problem 4. (Matrix Inequality, Matrix Norm) } (\textcolor{blue}{25 points})

\begin{enumerate}
    \item Let $\X \in \mathbb{R}^{n \times r}$ with $r \leq n$, and suppose that
    \[
    \|\X^\top \X - \I_r\|_2 = \tau < 1.
    \]
    Show that the smallest singular value of $\X$ satisfies
    \[
    \sigma_{\min}(\X) \geq 1 - \tau.
    \] (Hint: $||\A||_2=\max_{||\z||_2=1} ||\A\z||_2=\max_{||\z||_2=||\y||_2=1} |\y^T\A\z|$) (\textcolor{blue}{15 points})

    \item Assume matrix $E=ab^T$, where $a\in\mathbb{R}^m,b\in\mathbb{R}^n, a\neq\mathbf{0}_m,b\neq\mathbf{0}_n$, Prove the maximum singular value $\sigma_{1}(E)=||a||_2||b||_2$. (\textcolor{blue}{10 points})
    \end{enumerate}

\end{enumerate}

\noindent{\bf Solution:}








% =======================================================================================
\newpage
\noindent \textbf{Problem 5. (SVD)} (\textcolor{blue}{15 points})

Let $\theta_u$ and $\theta_v$ be real numbers that satisfy $0\le\theta_v-\theta_u\le \frac{\pi}{2}$ and
\[
A = \begin{bmatrix}
\cos(\theta_u) & \sin(\theta_u) \\
\cos(\theta_v) & \sin(\theta_v)
\end{bmatrix}.
\]
Find the singular value decomposition of $A$. (Hint: $\cos^2\left(\frac{\theta}{2}\right) = \frac{1+\cos(\theta)}{2},\sin^2\left(\frac{\theta}{2}\right) = \frac{1-\cos(\theta)}{2}$. The sum-to-product and product-to-sum identities are as follows.
% 和差化积（Sum-to-Product）
\[
\begin{aligned}
\sin A + \sin B &= 2\sin\!\left(\tfrac{A+B}{2}\right)\cos\!\left(\tfrac{A-B}{2}\right), \\
\sin A - \sin B &= 2\cos\!\left(\tfrac{A+B}{2}\right)\sin\!\left(\tfrac{A-B}{2}\right), \\
\cos A + \cos B &= 2\cos\!\left(\tfrac{A+B}{2}\right)\cos\!\left(\tfrac{A-B}{2}\right), \\
\cos A - \cos B &= -2\sin\!\left(\tfrac{A+B}{2}\right)\sin\!\left(\tfrac{A-B}{2}\right).
\end{aligned}
\]
% 积化和差（Product-to-Sum）
\[
\begin{aligned}
\sin\alpha\cos\beta &= \tfrac{1}{2}\big[\sin(\alpha+\beta) + \sin(\alpha-\beta)\big], \\
\cos\alpha\sin\beta &= \tfrac{1}{2}\big[\sin(\alpha+\beta) - \sin(\alpha-\beta)\big], \\
\cos\alpha\cos\beta &= \tfrac{1}{2}\big[\cos(\alpha+\beta) + \cos(\alpha-\beta)\big], \\
\sin\alpha\sin\beta &= -\tfrac{1}{2}\big[\cos(\alpha+\beta) - \cos(\alpha-\beta)\big].
\end{aligned}
\]
)


\bigskip
\noindent{\bf Solution:}




% =======================================================================================
\newpage
\noindent \textbf{Problem 6. (Sensitivity in Linear Systems)} \hfill (\textcolor{blue}{10 points})

Consider the least squares problem with a full column rank matrix $A\in \mathbb{R}^{m \times n}$ ($m \geq n$):
\[
\min_{x \in \mathbb{R}^n} \norm{Ax - y}_2^2.
\]
Let $\Delta y$ be a perturbation (disturbance) added to the observation vector $y$, and denote the resulting perturbation in the solution by $\Delta x$. Since $A$ has full column rank, the least-squares solution is unique and given by $x = A^\dagger y$, where $A^\dagger = (A^\top A)^{-1} A^\top$ is the Moore–Penrose pseudoinverse. Consequently,
\[
\Delta x = A^\dagger \, \Delta y.
\]

\begin{enumerate}
    \item Derive an inequality that relates the relative error in the solution to the relative error in the observation, i.e., prove $$
\frac{\norm{\Delta x}_2}{\norm{x}_2}
\leq  \kappa_2(A) \frac{\norm{\Delta y}_2}{\norm{y}_2}. $$
Where $\kappa_2(A)=\norm{A}_2 \, \norm{A^\dagger}_2 $. (Hint: Under the standard assumption in sensitivity analysis that the residual is small , we have $\norm{Ax}_2 \approx \norm{y}_2$, and thus
\[
\frac{1}{\norm{x}_2} \lesssim \frac{\norm{A}_2}{\norm{y}_2}.
\]) (\textcolor{blue}{7 points})
    \item Now let $A = \begin{bmatrix} A_1 \\ A_2 \end{bmatrix}$, where $A_1\in\mathbb{R}^{n\times n}$ is nonsingular and $A_2\in\mathbb{R}^{(m-n)\times n}$. Show that the upper bound of $||\Delta x||_2$ cannot infinitely increase with the increase of $m$. (\textit{Hint:}  $||\Delta x||_2=||A^{\dagger}\Delta y||_2$, and use the inequality $||A^\dagger||_2\le ||A_1^{-1}||_2$ to explain.) (\textcolor{blue}{3 points})
\end{enumerate}

\bigskip

\noindent \textbf{Solution:}


\end{document}
